\documentclass[a4paper]{article}

\usepackage[spanish]{babel}
\usepackage{graphicx}
\usepackage{amsmath, amssymb, mathrsfs}
\usepackage[margin=2cm]{geometry}
\usepackage{fancyhdr}
\usepackage{enumerate}
\usepackage[shortlabels]{enumitem}
\usepackage{parskip}
\usepackage[most]{tcolorbox}
\usepackage[hidelinks]{hyperref}
\usepackage{float}
\usepackage{booktabs}



% cabecera
\pagestyle{fancy}
\fancyhead[l]{Mecánica Celeste}
\fancyhead[c]{ProblemasV1}

\renewcommand{\headrulewidth}{0.2pt} % linea horizontal

\begin{document}

\begin{titlepage}
  \centering
  {\Large Universidad de Antioquia\par}
  \vspace{2cm}
  {\huge\bfseries Mecánica Celeste\par}
  \vspace{0.4cm}
  {\Large Problemas V1 \par}
  \vspace{2.5cm}
  {\large Autor:\par}
  \vspace{0.2cm}
  {\large Daniel Soto Villada\par}
  {\large Estudiante de Astronomía\par}
  \vfill
  {\large \today\par}
\end{titlepage}

\newpage
\tableofcontents
\newpage

\section{Problemas del libro de Harald}
\subsection{Ejemplo 4.4: El Hamiltoniano $H$ no siempre es la energía total $E$}

Sea $L = T - V$ el lagrangiano de un sistema de partículas con energía cinética $T$ y energía potencial $V$ independiente de la velocidad. Las ecuaciones que relacionan las coordenadas de posición $\vec{r}_i$ de las partículas con las $n$ coordenadas generalizadas $q_j$ dependen explícitamente del tiempo. Demuestre que
$$H = \sum_j \dot{q}_j p_j - L = T_1 - T_3 + V,$$
donde $T_1 = \sum_{k,l} A_{kl} \dot{q}_k \dot{q}_l$ y $T_3 = \sum_i \frac{1}{2} m_i \left(\frac{\partial \vec{r}_i}{\partial t}\right)^2$. ¿Qué se concluye si $\frac{\partial \vec{r}_i}{\partial t} = 0$?

\vspace{0.5cm}
\textbf{Solución:}

Sean las ecuaciones de transformación, para $i = 1, \dots, m$, dadas por $\vec{r}_i = \vec{r}_i(q_1, q_2, \dots, q_n, t)$. La derivada total respecto al tiempo $t$ proporciona la velocidad del $i$-ésimo punto de masa:
$$\vec{v}_i = \frac{d\vec{r}_i}{dt} = \sum_j \frac{\partial \vec{r}_i}{\partial q_j} \dot{q}_j + \frac{\partial \vec{r}_i}{\partial t}$$

Por lo tanto, la energía cinética $T$ del sistema es:
$$T = \frac{1}{2} \sum_i m_i \vec{v}_i \cdot \vec{v}_i = \frac{1}{2} \sum_i m_i \left( \sum_j \frac{\partial \vec{r}_i}{\partial q_j} \dot{q}_j + \frac{\partial \vec{r}_i}{\partial t} \right)^2$$

Al expandir el cuadrado, la energía cinética se puede descomponer en tres términos según el grado de las velocidades generalizadas $\dot{q}_j$:
$$T = \sum_{k,l} A_{kl} \dot{q}_k \dot{q}_l + \sum_k B_k \dot{q}_k + C = T_1 + T_2 + T_3$$
donde los coeficientes están definidos como:
$$A_{kl} = \frac{1}{2} \sum_i m_i \frac{\partial \vec{r}_i}{\partial q_k} \cdot \frac{\partial \vec{r}_i}{\partial q_l} = A_{lk}, \quad B_k = \sum_i m_i \frac{\partial \vec{r}_i}{\partial q_k} \cdot \frac{\partial \vec{r}_i}{\partial t}, \quad C = \frac{1}{2} \sum_i m_i \left(\frac{\partial \vec{r}_i}{\partial t}\right)^2$$

El Hamiltoniano $H$ se define mediante la transformación de Legendre como:
$$H = \sum_j \dot{q}_j p_j - L$$
donde el momento generalizado es $p_j = \frac{\partial L}{\partial \dot{q}_j}$. Dado que el potencial $V$ es independiente de las velocidades generalizadas $\dot{q}_j$, se cumple que $p_j = \frac{\partial T}{\partial \dot{q}_j}$. Sustituyendo esto y $L = T - V$ en la definición de $H$, obtenemos:
$$H = \sum_j \dot{q}_j \frac{\partial T}{\partial \dot{q}_j} - (T_1 + T_2 + T_3 - V)$$

Observamos que $T_1$ es una función homogénea de grado 2 en las velocidades generalizadas, $T_2$ es homogénea de grado 1 y $T_3$ es homogénea de grado 0. Aplicando el teorema de Euler para funciones homogéneas (Proposición 4.2 en Zuluaga), se tiene que:
$$\sum_j \dot{q}_j \frac{\partial T}{\partial \dot{q}_j} = 2T_1 + T_2$$

Sustituyendo este resultado en la expresión del Hamiltoniano:
$$H = (2T_1 + T_2) - (T_1 + T_2 + T_3 - V)$$
$$H = T_1 - T_3 + V$$

Si las ecuaciones de transformación son independientes del tiempo, es decir, si $\frac{\partial \vec{r}_i}{\partial t} = 0$ para todo $i$, entonces los términos $T_3$ y $T_2$ se anulan ($T_3 = 0$, $T_2 = 0$), y la energía cinética total se reduce a $T = T_1$. En este caso particular, la expresión del Hamiltoniano se simplifica a:
$$H = T_1 + V = T + V = E$$
donde $E$ representa la energía mecánica total del sistema.

Este ejemplo enfatiza un punto fundamental: el Hamiltoniano $H$ coincide con la energía mecánica total $E$ (en el caso de potenciales independientes de la velocidad) únicamente si las ecuaciones de transformación no involucran al tiempo $t$ explícitamente, es decir, cuando $\frac{\partial \vec{r}_i}{\partial t} = 0$ para todas las partículas del sistema.


\subsection{Ejemplo 4.8: El electrón en un magnetrón cilíndrico}

Un electrón de masa $m$ y carga $-e$ se mueve entre un alambre de radio $a$ con potencial eléctrico $\Phi = -\Phi_0$, y un conductor cilíndrico concéntrico conectado a tierra ($\Phi = 0$) de radio $R$. Existe un campo magnético constante $B$ paralelo al eje de la disposición. Deduzca la función Hamiltoniana y muestre que existen tres constantes de movimiento.


\textbf{Solución:}

Para describir el sistema, utilizamos coordenadas cilíndricas $(r, \theta, z)$. El potencial eléctrico $\Phi(r)$ que satisface las condiciones de frontera $\Phi(a) = -\Phi_0$ y $\Phi(R) = 0$ es:
$$ \Phi(r) = -\Phi_0 \frac{\ln(r/R)}{\ln(a/R)} $$
Por otro lado, el potencial vector magnético $\vec{A}$ asociado a un campo magnético uniforme $\vec{B} = B\hat{e}_z$ en el gauge simétrico es:
$$ \vec{A} = \frac{1}{2}Br\hat{e}_\theta $$
cuyas componentes son $(0, \frac{1}{2}Br, 0)$.

El Lagrangiano de una partícula cargada en un campo electromagnético está dado por $L = T - V + q\vec{v}\cdot\vec{A}$. Dado que la carga del electrón es $q = -e$, tenemos:
$$ L = \frac{1}{2}m(\dot{r}^2 + r^2\dot{\theta}^2 + \dot{z}^2) + e\Phi(r) - \frac{1}{2}eBr^2\dot{\theta} $$

Para pasar al formalismo hamiltoniano, primero debemos calcular los momentos canónicos conjugados $p_j = \frac{\partial L}{\partial \dot{q}_j}$ (Zuluaga, Sección 10.2):
$$ p_r = \frac{\partial L}{\partial \dot{r}} = m\dot{r} $$
$$ p_\theta = \frac{\partial L}{\partial \dot{\theta}} = mr^2\dot{\theta} - \frac{1}{2}eBr^2 $$
$$ p_z = \frac{\partial L}{\partial \dot{z}} = m\dot{z} $$

Despejamos las velocidades generalizadas en términos de los momentos canónicos:
$$ \dot{r} = \frac{p_r}{m} $$
$$ \dot{\theta} = \frac{1}{mr^2}\left(p_\theta + \frac{1}{2}eBr^2\right) $$
$$ \dot{z} = \frac{p_z}{m} $$

La función Hamiltoniana se define mediante la transformación de Legendre $H = \sum p_j \dot{q}_j - L$. Sustituyendo las expresiones anteriores:
$$ H = p_r\left(\frac{p_r}{m}\right) + p_\theta\left[\frac{1}{mr^2}\left(p_\theta + \frac{1}{2}eBr^2\right)\right] + p_z\left(\frac{p_z}{m}\right) - L $$
$$ H = \frac{p_r^2}{m} + \frac{p_\theta^2}{mr^2} + \frac{1}{2}eBp_\theta + \frac{p_z^2}{m} - \left[ \frac{1}{2}m\left(\frac{p_r^2}{m^2} + r^2\frac{(p_\theta + \frac{1}{2}eBr^2)^2}{m^2r^4} + \frac{p_z^2}{m^2}\right) + e\Phi(r) - \frac{1}{2}eBr^2\frac{p_\theta + \frac{1}{2}eBr^2}{mr^2} \right] $$
Simplificando algebraicamente los términos, obtenemos la función Hamiltoniana en términos de las variables canónicas:
$$ H(r, \theta, z, p_r, p_\theta, p_z) = \frac{p_r^2}{2m} + \frac{1}{2mr^2}\left(p_\theta + \frac{1}{2}eBr^2\right)^2 + \frac{p_z^2}{2m} - e\Phi(r) $$

Para identificar las constantes de movimiento, analizamos las simetrías del Hamiltoniano utilizando las ecuaciones canónicas de Hamilton (Zuluaga, Sección 10.3):

1. \textbf{Conservación de la energía:} Dado que el Lagrangiano y el Hamiltoniano no dependen explícitamente del tiempo ($\frac{\partial H}{\partial t} = 0$), por la \textbf{Proposición 10.2 (Conservación del Hamiltoniano)}, el propio Hamiltoniano $H$ es una constante de movimiento.

2. \textbf{Conservación del momento angular:} La coordenada angular $\theta$ no aparece explícitamente en el Hamiltoniano ($\frac{\partial H}{\partial \theta} = 0$). Por lo tanto, $\theta$ es una variable cíclica. De acuerdo con la \textbf{Proposición 10.1 (Variables cíclicas del Hamiltoniano)}, su momento canónico conjugado $p_\theta$ es una constante de movimiento.

3. \textbf{Conservación del momento lineal en $z$:} De manera análoga, la coordenada $z$ no aparece explícitamente en el Hamiltoniano ($\frac{\partial H}{\partial z} = 0$), por lo que es una variable cíclica. En consecuencia, su momento canónico conjugado $p_z$ también es una constante de movimiento.

En conclusión, las tres constantes de movimiento del sistema son el Hamiltoniano $H$, el momento canónico $p_\theta$ y el momento canónico $p_z$.


\subsection{Ejemplo 4.9: Un ejemplo con una cuantización problemática}


Intente resolver en la medida de lo posible el problema definido por el siguiente lagrangiano y obtenga el hamiltoniano correspondiente (el cual presenta un problema inherente al momento de ser cuantizado en mecánica cuántica):
$$
L(a, \phi, \dot{a}, \dot{\phi}) = -\frac{1}{2} \left[ a \dot{a}^2 + a^3 \dot{\phi}^2 + a^3 b \right] \qquad (4.98)
$$


\noindent \textbf{Solución:} \\
Usando la notación estándar para los momentos canónicos conjugados (definidos como $p_j = \frac{\partial L}{\partial \dot{q}_j}$, ver Sección 10.2 en Zuluaga), tenemos:
$$
p_a = \frac{\partial L}{\partial \dot{a}} = -a \dot{a}, \quad p_\phi = \frac{\partial L}{\partial \dot{\phi}} = -a^3 \dot{\phi} = \text{cte} \qquad (4.99)
$$
donde la constancia de $p_\phi$ se deduce del hecho de que $\frac{\partial L}{\partial \phi} = 0$ en la ecuación de Euler-Lagrange para $\phi$. Esto indica que $\phi$ es una \textit{variable cíclica} (o ignorable), y por el teorema de Noether, su momento generalizado conjugado se conserva (ver Definición 9.10 y Proposición 9.2 en Zuluaga).

La ecuación de Euler-Lagrange para la variable $a$ es:
$$
\frac{d}{dt} \left( \frac{\partial L}{\partial \dot{a}} \right) - \frac{\partial L}{\partial a} = 0 \implies \frac{d}{dt}(-a \dot{a}) - \left[ -\frac{1}{2} \dot{a}^2 - \frac{3}{2} a^2 \dot{\phi}^2 - \frac{3}{2} a^2 b \right] = 0 \qquad (4.100)
$$

Esta ecuación puede reescribirse de una forma más manejable reconociendo una derivada total (usando la regla de la cadena y la relación $d(a^2/2)/da = a$), lo que nos lleva a una cuadratura del sistema:
$$
\frac{d}{dt} \left[ \frac{\dot{a}^2}{2} + \frac{p_\phi^2}{2a^4} + \frac{b}{2} \right] = 0 \implies \frac{\dot{a}^2}{2} + \frac{p_\phi^2}{2a^4} + \frac{b}{2} = \text{cte} \qquad (4.101)
$$

De aquí se sigue que podemos separar variables para encontrar la evolución temporal de $a$:
$$
dt = \pm \frac{a^2 \, da}{\sqrt{p_\phi^2 + 2 \, \text{cte} \cdot a^4 - b a^4}} \qquad (4.102)
$$
La integral en el lado derecho es una integral elíptica y puede evaluarse con la ayuda de tablas de integrales elípticas o métodos numéricos adecuados.

Por otro lado, el hamiltoniano $H$ se obtiene mediante la transformación de Legendre del lagrangiano (ver Sección 10.3 en Zuluaga):
$$
H = p_\phi \dot{\phi} + p_a \dot{a} - L = -\frac{1}{2a} p_a^2 - \frac{1}{2a^3} p_\phi^2 - \frac{1}{2} a^3 b \qquad (4.103)
$$

Este tipo de hamiltoniano presenta un problema fundamental cuando se intenta cuantizar el sistema. La razón es que, en el hamiltoniano de la mecánica clásica, las cantidades canónicas $a$ y $p_a$ conmutan, es decir, pueden escribirse en cualquier orden sin alterar el resultado. Sin embargo, al usar este hamiltoniano clásico como punto de partida para su formulación en mecánica cuántica, se enfrenta el \textit{problema del ordenamiento de los operadores}. Es decir, surge la ambigüedad de cómo escribir el término que involucra $p_a$ y $a$, por ejemplo, como:
$$
\frac{1}{a} p_a \frac{1}{a} \quad \text{o bien} \quad p_a \frac{1}{a^2} \quad \text{o bien} \quad \frac{1}{a^2} p_a \qquad (4.104)
$$
ya que en mecánica cuántica los operadores $a$ y $p_a$ no conmutan ($[a, p_a] = i\hbar$). Afortunadamente, estos casos de ambigüedad de ordenamiento son raros en la práctica, y en las pocas ocasiones en que surgen, uno puede resolver la ambigüedad tomando, por ejemplo, un promedio simétrico (la mitad de una posibilidad y la mitad de la otra), o mediante criterios de hermiticidad del operador hamiltoniano.


\subsection{Ejemplo 7.7}

Un satélite tipo "mancuerna" (dumb-bell) consiste en dos masas puntuales iguales $m$ conectadas por una varilla rígida sin masa de longitud $2l$. Deduzca las ecuaciones de Euler-Lagrange del movimiento de este satélite en un plano bajo el campo gravitatorio de la Tierra (utilizando coordenadas polares $r$ y $\theta$ para el centro de masa, y el ángulo $\phi$ como el ángulo de inclinación del eje de la mancuerna respecto al vector de posición del centro de masa). ¿Cuál es la ley de conservación de la energía? Investigue en particular las siguientes soluciones especiales para $l/r_0 \ll 1$:
\begin{enumerate}
    \item[(a)] $r=r_0=\text{cte.}$, $\theta=\theta_0=\text{cte.}$, $\phi=0$ ("orientación de radio" o \textit{spoke orientation}).
    \item[(b)] $r=r_0=\text{cte.}$, $\theta=\theta_0=\text{cte.}$, $\phi=\pi/2$ ("orientación de lanza" o \textit{spear orientation}).
\end{enumerate}
¿En qué caso es la velocidad orbital más rápida o más lenta que la de una masa puntual? Finalmente, muestre que el movimiento de un satélite mancuerna suficientemente pequeño ($l/r$ pequeño) se subdivide en el movimiento de Kepler del centro de masa y un movimiento de rotación del satélite alrededor de su centro de masa.

\textbf{Solución:} \\

\textbf{1. Formulación del Lagrangiano} \\
La energía cinética $T$ del satélite es la suma de la energía cinética traslacional del centro de masa y la energía cinética de rotación alrededor del mismo. La velocidad del centro de masa en coordenadas polares es $v^2 = \dot{r}^2 + r^2\dot{\theta}^2$. La masa total es $2m$. El momento de inercia respecto al centro de masa es $I_{cm} = 2ml^2$, y su velocidad angular absoluta es $\dot{\theta} + \dot{\phi}$. Por lo tanto:
$$
T = \frac{1}{2}(2m)(\dot{r}^2 + r^2\dot{\theta}^2) + \frac{1}{2}(2ml^2)(\dot{\theta} + \dot{\phi})^2 = m(\dot{r}^2 + r^2\dot{\theta}^2) + ml^2(\dot{\theta} + \dot{\phi})^2
$$
La energía potencial gravitacional $V$ es la suma de las energías potenciales de las dos masas. Usando la ley de los cosenos para las distancias $r_1$ y $r_2$ desde el centro de la Tierra a cada masa:
$$
r_1^2 = r^2 + l^2 + 2rl\cos\phi, \quad r_2^2 = r^2 + l^2 - 2rl\cos\phi
$$
$$
V = -\frac{G M m}{r_1} - \frac{G M m}{r_2} = -m\mu \left[ (r^2+l^2+2rl\cos\phi)^{-1/2} + (r^2+l^2-2rl\cos\phi)^{-1/2} \right]
$$
donde $\mu = GM$ es el parámetro gravitacional de la Tierra (ver Sección 7.2 en Zuluaga). El Lagrangiano del sistema es $L = T - V$.

\textbf{2. Ecuaciones de Euler-Lagrange} \\
Aplicando las ecuaciones de Euler-Lagrange $\frac{d}{dt}\left(\frac{\partial L}{\partial \dot{q}_j}\right) - \frac{\partial L}{\partial q_j} = 0$ para las coordenadas generalizadas $q_j \in \{r, \theta, \phi\}$:

Para $r$:
$$
2m\ddot{r} - 2mr\dot{\theta}^2 - 2ml^2(\dot{\theta}+\dot{\phi})^2 - \frac{\partial V}{\partial r} = 0
$$
Dividiendo por $2m$ y evaluando la derivada del potencial:
$$
\ddot{r} - r\dot{\theta}^2 - l^2(\dot{\theta}+\dot{\phi})^2 = -\frac{\mu}{2} \left[ \frac{r+l\cos\phi}{(r^2+l^2+2rl\cos\phi)^{3/2}} + \frac{r-l\cos\phi}{(r^2+l^2-2rl\cos\phi)^{3/2}} \right]
$$

Para $\theta$:
$$
\frac{d}{dt} \left[ 2mr^2\dot{\theta} + 2ml^2(\dot{\theta}+\dot{\phi}) \right] = 0
$$
Esto refleja la conservación del momento angular total del sistema.

Para $\phi$:
$$
2ml^2(\ddot{\theta}+\ddot{\phi}) - \frac{\partial V}{\partial \phi} = 0
$$
Donde $\frac{\partial V}{\partial \phi} = m\mu r l \sin\phi \left[ (r^2+l^2+2rl\cos\phi)^{-3/2} - (r^2+l^2-2rl\cos\phi)^{-3/2} \right]$. Simplificando:
$$
l(\ddot{\theta}+\ddot{\phi}) + \mu r \sin\phi \left[ (r^2+l^2+2rl\cos\phi)^{-3/2} - (r^2+l^2-2rl\cos\phi)^{-3/2} \right] = 0
$$

\textbf{3. Conservación de la Energía} \\
Dado que el sistema es conservativo y el Lagrangiano no depende explícitamente del tiempo, la función de Jacobi coincide con la energía mecánica total $E = T + V$, la cual se conserva (Proposición 9.5 en Zuluaga):
$$
E = m(\dot{r}^2 + r^2\dot{\theta}^2) + ml^2(\dot{\theta} + \dot{\phi})^2 - m\mu \left[ (r^2+l^2+2rl\cos\phi)^{-1/2} + (r^2+l^2-2rl\cos\phi)^{-1/2} \right] = \text{cte.}
$$

\textbf{4. Soluciones Especiales para $l/r_0 \ll 1$} \\
Podemos expandir los términos de la fuerza usando la aproximación binomial para $l/r \ll 1$. Para $\phi$ constante ($\dot{\phi}=\ddot{\phi}=0$), la ecuación radial se aproxima a:
$$
\ddot{r} - r\dot{\theta}^2 \approx -\frac{\mu}{r^2} \left[ 1 + \frac{l^2}{r^2}(3\cos^2\phi - 1) \right]
$$

\begin{itemize}
    \item \textbf{(a) Orientación de radio ($\phi = 0$):} \\
    Para una órbita circular ($r=r_0$, $\ddot{r}=0$), la ecuación se convierte en:
    $$
    r_0\dot{\theta}^2 = \frac{\mu}{r_0^2} \left[ 1 + 2\left(\frac{l}{r_0}\right)^2 \right] \implies \dot{\theta}^2 = \frac{\mu}{r_0^3} \left[ 1 + 2\left(\frac{l}{r_0}\right)^2 \right] > \frac{\mu}{r_0^3}
    $$
    La velocidad orbital es \textbf{más rápida} que la de una masa puntual en la misma órbita.

    \item \textbf{(b) Orientación de lanza ($\phi = \pi/2$):} \\
    De manera similar, para $\phi = \pi/2$ y $\ddot{r}=0$:
    $$
    r_0\dot{\theta}^2 = \frac{\mu}{r_0^2} \left[ 1 - \left(\frac{l}{r_0}\right)^2 \right] \implies \dot{\theta}^2 = \frac{\mu}{r_0^3} \left[ 1 - \left(\frac{l}{r_0}\right)^2 \right] < \frac{\mu}{r_0^3}
    $$
    La velocidad orbital es \textbf{más lenta} que la de una masa puntual.
\end{itemize}

\textbf{5. Desacoplamiento del Movimiento} \\
Expandiendo el potencial a primer orden en $(l/r)^2$, obtenemos:
$$
V \approx -\frac{2m\mu}{r} \left[ 1 + \frac{l^2}{2r^2}(3\cos^2\phi - 1) \right]
$$
Sustituyendo esto en las ecuaciones de movimiento para $r$ y $\theta$, los términos que contienen $l^2/r^2$ son de orden superior y pueden despreciarse en primera aproximación. Las ecuaciones para el centro de masa se reducen a:
$$
\ddot{r} - r\dot{\theta}^2 \approx -\frac{\mu}{r^2}, \quad \frac{d}{dt}(r^2\dot{\theta}) \approx 0
$$
Estas son exactamente las ecuaciones del problema de Kepler para el centro de masa (ver Sección 7.2 en Zuluaga). 

Por otro lado, la ecuación para $\phi$ se convierte en:
$$
l^2\ddot{\phi} + l^2\ddot{\theta} + \frac{3\mu l^2}{2r^3} \sin(2\phi) \approx 0
$$
Dado que $\ddot{\theta} \sim \mathcal{O}(1/r^3)$, el término $l^2\ddot{\theta}$ es del mismo orden que el término de torque gravitatorio (efecto de marea). Así, la rotación $\phi$ se describe por una ecuación diferencial que depende paramétricamente de la solución de Kepler $r(t)$, demostrando que el movimiento se subdivide en el movimiento traslacional de Kepler del centro de masa y un movimiento de rotación independiente (acoplado solo a través de las fuerzas de marea) alrededor de dicho centro de masa.

\subsection{Ejemplo 7.8}

Un satélite que orbita la Tierra en una órbita circular de radio $r_1$ debe ser elevado a una órbita circular concéntrica de radio $r_2 > r_1$ mediante dos impulsos de cohete. La cantidad requerida de combustible se optimiza si, como se indica en la figura correspondiente, la transición se realiza a través de una órbita elíptica con la Tierra en uno de sus focos (conocida como órbita de transferencia de Hohmann). Calcule, en función de $r_1$ y $r_2$, el incremento de velocidad requerido para la transición, es decir, la cantidad:
$$
\Delta v = (v_{\max} - v_{r1}) + (v_2 - v_{\min})
$$
donde $v_{r1}$ y $v_2$ son las velocidades en las órbitas circulares inicial y final, respectivamente, y $v_{\max}$ y $v_{\min}$ son las velocidades máxima y mínima en la órbita elíptica de transferencia.

\textbf{Solución:} \\
Las velocidades de las órbitas circulares se obtienen directamente de la mecánica orbital (o de la ecuación de la vis viva con excentricidad $e=0$):
$$
v_{r1} = \sqrt{\frac{\mu}{r_1}}, \quad v_2 = \sqrt{\frac{\mu}{r_2}}
$$
donde $\mu = GM$ es el parámetro gravitacional de la Tierra (ver Sección 7.2 en Zuluaga).

Para la órbita elíptica de transferencia, la velocidad en cualquier punto está dada por la ecuación de la \textit{vis viva}:
$$
v^2 = \mu \left( \frac{2}{r} - \frac{1}{a} \right)
$$
donde $a$ es el semieje mayor de la elipse. Dado que la elipse es tangente a las órbitas circulares en el periapsis y el apoapsis, tenemos que las distancias extremales son:
$$
r_{\min} = r_1 = a(1-e) \quad \text{y} \quad r_{\max} = r_2 = a(1+e)
$$
Sumando estas dos ecuaciones, obtenemos el semieje mayor de la órbita de transferencia:
$$
r_1 + r_2 = 2a \implies a = \frac{r_1 + r_2}{2}
$$
La velocidad máxima $v_{\max}$ ocurre en el periapsis de la elipse ($r = r_1$). Sustituyendo $a$ en la ecuación de la vis viva:
$$
v_{\max} = \sqrt{\mu \left( \frac{2}{r_1} - \frac{2}{r_1 + r_2} \right)} = \sqrt{\frac{2\mu r_2}{(r_1 + r_2)r_1}}
$$
De manera análoga, la velocidad mínima $v_{\min}$ ocurre en el apoapsis de la elipse ($r = r_2$):
$$
v_{\min} = \sqrt{\mu \left( \frac{2}{r_2} - \frac{2}{r_1 + r_2} \right)} = \sqrt{\frac{2\mu r_1}{(r_1 + r_2)r_2}}
$$
Finalmente, sustituimos estas expresiones en la fórmula del incremento total de velocidad $\Delta v$ planteada en el enunciado:
$$
\Delta v = \left( \sqrt{\frac{2\mu r_2}{(r_1 + r_2)r_1}} - \sqrt{\frac{\mu}{r_1}} \right) + \left( \sqrt{\frac{\mu}{r_2}} - \sqrt{\frac{2\mu r_1}{(r_1 + r_2)r_2}} \right)
$$
Factorizando el término $\sqrt{\mu}$, la expresión final y compacta para el incremento de velocidad requerido es:
$$
\Delta v = \sqrt{\mu} \left( \sqrt{\frac{2r_2}{r_1(r_1 + r_2)}} - \frac{1}{\sqrt{r_1}} + \frac{1}{\sqrt{r_2}} - \sqrt{\frac{2r_1}{r_2(r_1 + r_2)}} \right)
$$
Este resultado demuestra cómo el $\Delta v$ total depende exclusivamente de los radios de las órbitas inicial y final, y del parámetro gravitacional del cuerpo central.

\bibliographystyle{plainurl}
\bibliography{bibliografia} 
\end{document}
